\section{Impact Statement}

This work exposes a critical security risk in LLM-assisted code generation: 
models can embed malicious URLs in response to developer-style prompts. 

\noindent \textbf{Advancing Model Safety.} 
The primary impact of this work is improved awareness and detection 
tooling, enabling researchers and practitioners to audit production LLMs and 
design effective mitigations. 
The released evaluation resources facilitate reproducible 
security testing for new and existing LLMs.

\noindent \textbf{Ethical Considerations.} 
We acknowledge that publicizing failure patterns could theoretically lower 
the barrier for crafting adversarial inputs. 
We mitigate this risk by focusing on measurement methodologies 
rather than exploit instructions and by actively reporting 
discovered scam domains to database maintainers. 
Given that this threat vector is already active in the wild 
(Section~\ref{sec:example}), this work serves to demystify an 
existing security gap rather than introducing a novel threat.



% This work analyzes a security risk in LLM-assisted code generation: 
% models can generate code that embeds malicious URLs in response to developer-style prompts. 
% The primary positive impact is improved awareness and diagnostic tooling for this failure mode, 
% which can help researchers and practitioners evaluate model safety, 
% audit training data pipelines, and design mitigations. 
% The released evaluation resources can support reproducible security testing new and existing LLMs.
% There are potential negative impacts. Publicizing concrete failure patterns and releasing 
% prompts could lower the barrier for attackers to probe systems or craft inputs 
% that elicit unsafe outputs. We mitigate this by focusing on auditing and measurement 
% rather than providing exploit instructions, and by reporting newly discovered scam domains to database maintainers.
% Overall, we believe the benefits of exposing and measuring 
% this security risk outweigh the risks, provided that responsible disclosure and defensive practices are 
% followed. Given that this threat vector is already active in the wild 
% (Section~\ref{sec:example}), this work serves to demystify an existing 
% security gap rather than introducing a novel threat.
